\documentclass[11pt]{amsart}
\usepackage{amsmath,amssymb,amsthm,mathtools}
\usepackage[margin=1in]{geometry}
\usepackage{enumerate}
\usepackage{hyperref}

\newtheorem{theorem}{Theorem}[section]
\newtheorem*{theorem*}{Theorem}
\newtheorem{proposition}[theorem]{Proposition}
\newtheorem{lemma}[theorem]{Lemma}
\newtheorem{corollary}[theorem]{Corollary}
\theoremstyle{definition}
\newtheorem{definition}[theorem]{Definition}
\newtheorem{problem}[theorem]{Problem}
\theoremstyle{remark}
\newtheorem{remark}[theorem]{Remark}
\newtheorem{observation}[theorem]{Observation}

\newcommand{\bbR}{\mathbb{R}}
\newcommand{\bbC}{\mathbb{C}}
\DeclareMathOperator{\Real}{Re}

\title[$\zeta$ is not Stepanov a.p.]{The Riemann zeta function is not Stepanov almost periodic\\
in any sub-strip of the critical strip}

\author{David Alejandro Trejo Pizzo}
\thanks{Independent Researcher. \texttt{dtrejopizzo@gmail.com}}


\begin{document}

\begin{abstract}
For $\sigma>1$ the Riemann zeta function is Bohr almost periodic on
vertical lines, and inside the critical strip it is almost periodic in
the Besicovitch $B^2$ sense. Between the uniform and the Besicovitch
scales sits the Stepanov scale, the only one of the classical
almost-periodicity scales whose seminorm sees windows of fixed length.
We prove, unconditionally, that $\zeta$ is not Stepanov almost periodic,
for any exponent $p\geq1$, on any closed sub-strip
$[\sigma_1,\sigma_2]\subset(\tfrac12,1)$, in the band (vector-valued)
sense: the vectorial Stepanov norm
$\sup_x\bigl(\int_x^{x+1}\int_{\sigma_1}^{\sigma_2}
|\zeta(\sigma+it)|^p\,d\sigma\,dt\bigr)^{1/p}$ is infinite, so
$t\mapsto\zeta(\cdot+it)$, viewed as an
$L^p(\sigma_1,\sigma_2)$-valued function, is not $S^p$-almost periodic.
The proof combines the unconditional $\Omega$-theorems for large values
of $\zeta$ on fixed interior lines (Titchmarsh; Aistleitner) with the
subharmonicity of $|\zeta|^p$, which promotes a large point value to a
large area integral over a band-times-unit-window rectangle. The
fixed-line scalar version of the statement is not claimed and remains
open; we explain precisely why the method does not decide it. The
result places $\zeta$ inside a clean dichotomy: the seminorms under
which $\zeta$ \emph{is} almost periodic in the strip ($B^2$, Weyl)
are blind to compact windows, while the seminorms that see compact
windows (uniform, Stepanov) are destroyed by the growth of $\zeta$.
\end{abstract}

\maketitle

\tableofcontents
\newpage

\section{Introduction}\label{sec:intro}

\subsection{Almost periodicity of $\zeta$: the two classical scales}
The connection between the Riemann zeta function and almost periodicity
is as old as the theory of almost periodic functions itself: Bohr
created the theory in large part as a tool for the Riemann Hypothesis
\cite{Bohr22}. Two facts anchor the subject.

First, in any closed half-plane $\sigma\geq1+\delta$ ($\delta>0$) the
Dirichlet series $\sum n^{-s}$ converges absolutely and uniformly, so
$t\mapsto\zeta(\sigma+it)$ is a uniform limit of trigonometric
polynomials $\sum_{n\leq N}n^{-\sigma}e^{-it\log n}$ and is therefore
Bohr (uniformly) almost periodic, on each line and indeed uniformly in
the half-plane \cite{Bohr22,Besicovitch}. This fails on every line
$\sigma\leq1$, simply because $t\mapsto\zeta(\sigma+it)$ is unbounded
there (see the $\Omega$-theorems quoted in \S\ref{sec:main}), and Bohr
almost periodic functions are bounded.

Second, inside the critical strip the almost periodicity survives in
an averaged sense. For fixed $\sigma>\tfrac12$ the partial sums of the
Dirichlet series still approximate $\zeta$ in quadratic mean over long
intervals: by the classical mean-value theorems
(Hardy--Littlewood; see \cite[Ch.~VII]{Titchmarsh}),
\[
\limsup_{T\to\infty}\frac{1}{2T}\int_{-T}^{T}
\Bigl|\zeta(\sigma+it)-\sum_{n\leq N}n^{-\sigma-it}\Bigr|^2 dt
\;\xrightarrow[N\to\infty]{}\;0 ,
\]
so $t\mapsto\zeta(\sigma+it)$ is almost periodic in the Besicovitch
$B^2$ sense \cite[Ch.~II]{Besicovitch} on every line
$\sigma>\tfrac12$. Thus $\zeta$ is almost periodic in the strip---but
only for a seminorm that averages over windows of \emph{growing}
length.

\subsection{The intermediate scale, and why it matters}
Between the uniform norm and the Besicovitch seminorm sits the
classical hierarchy of almost periodicity classes
\cite[Ch.~II]{Besicovitch}, \cite{LevitanZhikov}:
\[
\text{Bohr}\ \subset\ S^p\ \subset\ W^p\ \subset\ B^p ,
\]
where $S^p$ is the Stepanov class, with seminorm
$\|f\|_{S^p}=\sup_x\bigl(\int_x^{x+1}|f|^p\bigr)^{1/p}$, and $W^p$,
$B^p$ are the Weyl and Besicovitch classes, whose seminorms average
over windows of length $L\to\infty$. The Stepanov seminorm is the only
intermediate one that \emph{sees windows of fixed length}: it controls
the $L^p$ mass of $f$ on every interval of length $1$, uniformly in
the position of the interval. The Weyl and Besicovitch seminorms do
not: for any $f\in L^p_{\mathrm{loc}}$ and any compact $K$, the Weyl
and Besicovitch seminorms of $f\cdot\mathbf{1}_K$ vanish, because the
average of a fixed compact contribution over windows of length
$L\to\infty$ decays like $1/L$, uniformly in the window position.

This makes the Stepanov scale the natural question for $\zeta$ in the
strip. Any almost-periodicity statement strong enough to transport
local information---values, or zeros via Rouch\'e-type arguments---from
one window of the critical strip to another must control fixed
windows, which is exactly what $B^2$ almost periodicity does not do
and Stepanov almost periodicity would. The question is therefore:
\emph{does the almost periodicity of $\zeta$ in the strip survive at
the Stepanov scale?}

\subsection{The result}
The answer is no, unconditionally, in the band (vector-valued) sense,
and the reason is growth: the same unboundedness that kills Bohr
almost periodicity on the lines $\sigma\leq1$ kills the Stepanov
seminorm in the strip, once the pointwise large values are converted
into area integrals by subharmonicity.

\begin{theorem*}[Main theorem; see Theorem~\ref{thm:main}]
Let $B=[\sigma_1,\sigma_2]\subset(\tfrac12,1)$ with
$\sigma_1<\sigma_2$, and let $p\geq1$. Then the vectorial Stepanov norm
\[
\|\zeta\|_{S^p(B)}
:=\sup_{x\in\bbR}\Bigl(\int_x^{x+1}\!\!\int_{\sigma_1}^{\sigma_2}
|\zeta(\sigma+it)|^p\,d\sigma\,dt\Bigr)^{1/p}
\]
is infinite. Consequently $t\mapsto\zeta(\cdot+it)$, viewed as an
$L^p(\sigma_1,\sigma_2)$-valued function of $t\in\bbR$, is not
$S^p$-almost periodic, for any sub-strip $B$ of the open critical
strip and any $p\geq1$.
\end{theorem*}

The statement and proof concern the \emph{band} norm: $\zeta$ is
treated as a function of the single real variable $t$ with values in
the Banach space $L^p(\sigma_1,\sigma_2)$, and the Stepanov seminorm
is the classical scalar one with $|f(t)|$ replaced by
$\|\zeta(\cdot+it)\|_{L^p(\sigma_1,\sigma_2)}$. This is the standard
vector-valued (Banach-space-valued) almost periodicity of
Levitan--Zhikov \cite{LevitanZhikov}. We emphasise, and discuss in
\S\ref{sec:discussion}, that the fixed-line scalar version---whether
$t\mapsto\zeta(\sigma_0+it)$ can be $S^p$-almost periodic for some
single $\sigma_0\in(\tfrac12,1)$---is \emph{not} decided by our
method and remains open.

The proof is short and entirely classical in its ingredients: an
unconditional $\Omega$-theorem on a fixed interior line
(Titchmarsh \cite[Ch.~VIII]{Titchmarsh}; quantitatively Aistleitner
\cite{Aistleitner}), the sub-mean-value inequality for the subharmonic
function $|\zeta|^p$ on disks of a fixed radius contained in the band,
and a monotone inclusion of domains. What may be of interest beyond
the statement itself is the position the result occupies: together
with the classical positive results it exhibits an exact dichotomy in
seminorm coordinates (Observation~\ref{obs:dichotomy})---every
almost-periodicity seminorm that sees compact windows is destroyed by
the growth of $\zeta$ in the strip, and every seminorm under which
$\zeta$ is almost periodic there is blind to compact windows.

\section{Stepanov almost periodicity in the band sense}\label{sec:defs}

Throughout, $B=[\sigma_1,\sigma_2]\subset(\tfrac12,1)$ is a closed
band with $\sigma_1<\sigma_2$, and $p\geq1$.

\begin{definition}[Vectorial Stepanov norm]\label{def:norm}
For a function $F\colon\bbR\to L^p(\sigma_1,\sigma_2)$, locally
$p$-integrable, set
\[
\|F\|_{S^p(B)}^p:=\sup_{x\in\bbR}\int_x^{x+1}
\|F(t)\|_{L^p(\sigma_1,\sigma_2)}^p\,dt .
\]
For $F(t)=\zeta(\cdot+it)$ this is
\[
\|\zeta\|_{S^p(B)}^p
=\sup_{x\in\bbR}\int_x^{x+1}\!\!\int_{\sigma_1}^{\sigma_2}
|\zeta(\sigma+it)|^p\,d\sigma\,dt .
\]
\end{definition}

Since $\zeta$ is holomorphic on the closed band $B\times\bbR$ (the
pole $s=1$ lies outside $B$), the map $t\mapsto\zeta(\cdot+it)$ is
continuous from $\bbR$ into $L^p(\sigma_1,\sigma_2)$ ($\zeta$ is
uniformly continuous on $B\times K$ for every compact
$K\subset\bbR$), so all the integrals above are finite for each fixed
$x$ and the supremum is well defined in $[0,\infty]$.

\begin{definition}[$S^p$-almost periodicity, band sense]\label{def:sp-ap}
A locally $p$-integrable $F\colon\bbR\to L^p(\sigma_1,\sigma_2)$ is
\emph{$S^p$-almost periodic} if either of the following equivalent
conditions holds:
\begin{enumerate}[(i)]
\item (closure of trigonometric polynomials) for every
$\varepsilon>0$ there is a trigonometric polynomial
$P(t)=\sum_{k=1}^{N}a_k e^{i\lambda_k t}$ with coefficients
$a_k\in L^p(\sigma_1,\sigma_2)$ and real frequencies $\lambda_k$ such
that $\|F-P\|_{S^p(B)}<\varepsilon$;
\item (relatively dense almost periods) for every $\varepsilon>0$ the
set of $\varepsilon$-almost periods
$\{\tau\in\bbR:\ \|F(\cdot+\tau)-F(\cdot)\|_{S^p(B)}<\varepsilon\}$
is relatively dense in $\bbR$, i.e.\ meets every interval of some
fixed length $L(\varepsilon)$.
\end{enumerate}
\end{definition}

The equivalence of (i) and (ii) is the classical approximation theorem
of the Stepanov theory, due in the scalar case to Besicovitch and
others \cite[Ch.~II]{Besicovitch} and valid verbatim for functions
with values in a Banach space \cite{LevitanZhikov}; equivalently, the
Stepanov theory is the Bohr theory applied to the Bochner transform
$t\mapsto F(t+\cdot)\in L^p\bigl((0,1);L^p(\sigma_1,\sigma_2)\bigr)$.
We do not reprove the equivalence; the only consequence of
almost periodicity we use is the following finiteness lemma, which we
prove under \emph{both} definitions so that the main theorem is
independent of which one the reader prefers.

\begin{lemma}[Almost periodicity implies a finite Stepanov norm]
\label{lem:bounded}
If $F$ is $S^p$-almost periodic in the sense of
Definition~\ref{def:sp-ap}, and $t\mapsto\|F(t)\|_{L^p}$ is locally
bounded, then $\|F\|_{S^p(B)}<\infty$.
\end{lemma}

\begin{proof}
\emph{Via (i).} Choose $P$ with $\|F-P\|_{S^p(B)}<1$. A trigonometric
polynomial is a bounded function of $t$ with values in
$L^p(\sigma_1,\sigma_2)$: $\|P(t)\|_{L^p}\leq\sum_k\|a_k\|_{L^p}=:M$
for all $t$, whence $\|P\|_{S^p(B)}\leq M<\infty$. By the triangle
inequality for the seminorm $\|\cdot\|_{S^p(B)}$,
\[
\|F\|_{S^p(B)}\;\leq\;\|F-P\|_{S^p(B)}+\|P\|_{S^p(B)}\;<\;1+M\;<\;\infty .
\]

\emph{Via (ii).} Take $\varepsilon=1$ and let $L=L(1)$ be such that
every interval of length $L$ contains a $1$-almost period. Given
$x\in\bbR$, pick a $1$-almost period $\tau\in[x,x+L]$ and set
$y:=x-\tau\in[-L,0]$. Then
\[
\Bigl(\int_x^{x+1}\|F(t)\|_{L^p}^p\,dt\Bigr)^{1/p}
=\Bigl(\int_y^{y+1}\|F(t+\tau)\|_{L^p}^p\,dt\Bigr)^{1/p}
\leq 1+\Bigl(\int_y^{y+1}\|F(t)\|_{L^p}^p\,dt\Bigr)^{1/p},
\]
and the last quantity is bounded by
$1+\sup_{y\in[-L,0]}\bigl(\int_y^{y+1}\|F\|_{L^p}^p\bigr)^{1/p}$,
which is finite by local boundedness. The bound is uniform in $x$.
\end{proof}

Note that the conclusion is finiteness of the \emph{norm}
$\|F\|_{S^p(B)}$, not pointwise boundedness of $F$: an $S^p$ limit of
bounded functions need not be pointwise bounded, but its $S^p$ norm is
finite, and that is all the main theorem contradicts. For
$F(t)=\zeta(\cdot+it)$ the local boundedness hypothesis holds, since
$\zeta$ is continuous on the closed band.

\section{The main theorem}\label{sec:main}

The analytic input is an unconditional lower bound for large values of
$\zeta$ on a fixed line in the interior of the strip. Two classical
sources suffice, and either one alone is enough for our purposes.

\begin{theorem}[Titchmarsh {\cite[Ch.~VIII]{Titchmarsh}}]
\label{thm:titchmarsh}
For fixed $\sigma\in(\tfrac12,1)$ and every $\varepsilon>0$,
\[
|\zeta(\sigma+it)|=\Omega\Bigl(\exp\bigl((\log t)^{1-\sigma-\varepsilon}\bigr)\Bigr)
\qquad(t\to\infty),
\]
unconditionally. In particular $t\mapsto\zeta(\sigma+it)$ is unbounded
on every such line.
\end{theorem}

\begin{theorem}[Aistleitner {\cite{Aistleitner}}]\label{thm:aistleitner}
For fixed $\alpha\in(\tfrac12,1)$ there is a constant
$c_\alpha=0.18\,(2\alpha-1)^{1-\alpha}$ such that, unconditionally,
for all sufficiently large $T$,
\[
\max_{0\leq t\leq T}|\zeta(\alpha+it)|
\;\geq\;\exp\Bigl(c_\alpha\,\frac{(\log T)^{1-\alpha}}{(\log\log T)^{\alpha}}\Bigr).
\]
\end{theorem}

Theorem~\ref{thm:aistleitner} sharpens, by the resonance method, a
result going back to Montgomery \cite{Montgomery}; we stress that the
unconditional statements we rely on are
Theorems~\ref{thm:titchmarsh} and~\ref{thm:aistleitner} as stated, and
that the quantitative rates play no role below: the proof needs only
the qualitative conclusion that $\zeta$ is unbounded on one interior
line of each sub-band.

\begin{theorem}[Main theorem]\label{thm:main}
Let $B=[\sigma_1,\sigma_2]\subset(\tfrac12,1)$ with
$\sigma_1<\sigma_2$, and let $p\geq1$. Then
\[
\|\zeta\|_{S^p(B)}=\infty .
\]
In particular, by Lemma~\ref{lem:bounded}, $t\mapsto\zeta(\cdot+it)$
is not $S^p$-almost periodic as an $L^p(\sigma_1,\sigma_2)$-valued
function, under either definition in Definition~\ref{def:sp-ap}.
\end{theorem}

\begin{proof}
Fix $\sigma\in(\sigma_1,\sigma_2)$ and set
\[
r:=\tfrac12\min\bigl(\sigma-\sigma_1,\ \sigma_2-\sigma,\ \tfrac14\bigr)>0 ,
\]
so that $r\leq\tfrac18$ and every disk $D(\sigma+it,r)$ with
$t\in\bbR$ satisfies the inclusion
\begin{equation}\label{eq:inclusion}
D(\sigma+it,\,r)\ \subset\ [\sigma_1,\sigma_2]\times[t-r,\,t+r],
\end{equation}
because $r\leq(\sigma-\sigma_1)/2$ and $r\leq(\sigma_2-\sigma)/2$.

\medskip
\emph{Step 1: a sequence of large values.}
By Theorem~\ref{thm:titchmarsh} (or Theorem~\ref{thm:aistleitner})
there exists a sequence $t_n\to\infty$ with
\[
V_n:=|\zeta(\sigma+it_n)|\longrightarrow\infty .
\]
For an $\Omega$-statement this is immediate (a sequence realising the
$\limsup$). For a statement of the form
$\max_{0\leq t\leq T}|\zeta(\sigma+it)|\geq\Phi(T)\to\infty$, let
$t_T$ be a maximiser on $[0,T]$; the $t_T$ cannot remain in a compact
set, since $\zeta$ is continuous, hence bounded, on
$[\sigma_1,\sigma_2]\times[0,K]$ for every $K$, while
$|\zeta(\sigma+it_T)|\geq\Phi(T)\to\infty$; so $t_T\to\infty$ along a
subsequence, which we take as $(t_n)$. We may also assume $t_n>1$ for
all $n$, so that the disks below stay away from the real axis.

\medskip
\emph{Step 2: subharmonicity.}
The function $|\zeta|^p$ is subharmonic on a neighbourhood of
$B\times[1,\infty)$: $\zeta$ is holomorphic there, $\log|\zeta|$ is
subharmonic, and $|\zeta|^p=\exp(p\log|\zeta|)$ is the composition of
a subharmonic function with the convex increasing function
$\exp(p\,\cdot)$, hence subharmonic (this holds for every $p>0$; the
restriction $p\geq1$ only fixes the classical Stepanov class). The
sub-mean-value inequality on the disk $D(\sigma+it_n,r)$ gives
\[
V_n^p\;\leq\;\frac{1}{\pi r^2}\iint_{D(\sigma+it_n,\,r)}
|\zeta|^p\,dA .
\]

\medskip
\emph{Step 3: inclusion of domains.}
The integrand is nonnegative, so by \eqref{eq:inclusion},
\[
V_n^p\;\leq\;\frac{1}{\pi r^2}
\int_{t_n-r}^{t_n+r}\!\!\int_{\sigma_1}^{\sigma_2}
|\zeta(\sigma'+it)|^p\,d\sigma'\,dt .
\]

\medskip
\emph{Step 4: the window fits.}
Since $2r\leq\tfrac14<1$, the time interval $[t_n-r,t_n+r]$ is
contained in the unit window $[x_n,x_n+1]$ with $x_n:=t_n-r$. Hence
\[
\|\zeta\|_{S^p(B)}^p\;\geq\;
\int_{x_n}^{x_n+1}\!\!\int_{\sigma_1}^{\sigma_2}
|\zeta(\sigma'+it)|^p\,d\sigma'\,dt
\;\geq\;\pi r^2\,V_n^p\;\longrightarrow\;\infty ,
\]
and the supremum over $x$ is infinite. Each individual window integral
is finite ($\zeta$ is continuous on $B\times$ compacts); the
divergence is genuinely a divergence of the supremum along
$x_n\to\infty$, as it must be.
\end{proof}

\begin{remark}[What the proof does and does not use]\label{rem:robust}
(a) The argument never converts the two-dimensional sub-mean value
into a one-dimensional average at a fixed abscissa: the quantity shown
to be infinite \emph{is} the two-dimensional band-times-window
average, and Step~3 is a monotone inclusion of domains with a
nonnegative integrand---no Fubini, no choice of a distinguished
vertical line inside the disk. This is precisely why the theorem is a
statement about the band norm and not about any fixed line; see
Problem~\ref{prob:fixedline}.
(b) The growth \emph{rate} in
Theorems~\ref{thm:titchmarsh}--\ref{thm:aistleitner} is irrelevant:
$S^p$-almost periodicity would impose one fixed finite bound
$\|\zeta\|_{S^p(B)}=C<\infty$, and any sequence $V_n\to\infty$
violates $\pi r^2 V_n^p\leq C^p$. Mere unboundedness of $\zeta$ on a
single interior line of the sub-band suffices.
(c) The proof works for every $p>0$ and every sub-band, with the
radius $r$ calibrated locally; the hypothesis $\sigma_1<\sigma_2$ is
essential (for $\sigma_1=\sigma_2$ there is no band and no disk).
\end{remark}

\begin{remark}[Non-membership, not just non-almost-periodicity]
The theorem proves more than the failure of almost periodicity:
$\|\zeta\|_{S^p(B)}=\infty$ means $\zeta$ does not even belong to the
Stepanov space on which the class of $S^p$-almost periodic functions
is modelled. The prose statement ``$\zeta$ is not $S^p$-almost
periodic on any sub-strip, in the band sense'' is the correct
consequence via Lemma~\ref{lem:bounded}.
\end{remark}

\section{Discussion}\label{sec:discussion}

\subsection{A dichotomy in seminorm coordinates}

\begin{observation}[Seeing windows versus surviving growth]
\label{obs:dichotomy}
In the strip $\tfrac12<\sigma<1$, the classical almost-periodicity
seminorms split into two groups, and $\zeta$ separates them exactly.
\begin{enumerate}[(a)]
\item \emph{Seminorms blind to compact windows ($B^p$, $W^p$):} the
Besicovitch and Weyl seminorms average over windows of length
$L\to\infty$, so the seminorm of $f\cdot\mathbf{1}_K$ vanishes for
every compact $K$ (for Weyl, note that the supremum over the window
position is taken \emph{before} the limit $L\to\infty$, and the
average of a fixed compact contribution still decays like $1/L$
uniformly in the position). Under these seminorms $\zeta$ \emph{is}
almost periodic in the strip ($B^2$, classically, from the mean-square
approximation by partial sums; \S\ref{sec:intro}).
\item \emph{Seminorms that see compact windows (uniform, $S^p$):} the
uniform norm and the Stepanov seminorms control fixed windows
uniformly in their position. Under these, almost periodicity of
$\zeta$ fails in the strip: uniformly because $\zeta$ is unbounded on
every line $\sigma<1$ (Theorem~\ref{thm:titchmarsh}), and in the
Stepanov band sense by Theorem~\ref{thm:main}.
\end{enumerate}
In short: \emph{every seminorm under which $\zeta$ is almost periodic
in the strip is blind to compact windows, and every seminorm that sees
compact windows is destroyed by the growth of $\zeta$.} The mechanism
of failure is in both cases growth, not arithmetic: the unconditional
large values of $\zeta$ are incompatible with any almost-periodicity
seminorm that assigns mass to windows of fixed length.
\end{observation}

\subsection{The open fixed-line problem}

\begin{problem}[Fixed-line scalar Stepanov almost periodicity]
\label{prob:fixedline}
Is $t\mapsto\zeta(\sigma_0+it)$ $S^p$-almost periodic, as a scalar
function, for some $\sigma_0\in(\tfrac12,1)$ and some $p\geq1$?
Equivalently, in the bounded direction: is
$\sup_x\int_x^{x+1}|\zeta(\sigma_0+it)|^p\,dt$ finite for some such
$\sigma_0$?
\end{problem}

We claim nothing about this problem; in particular,
Theorem~\ref{thm:main} does \emph{not} decide it, and the reason is
structural. The sub-mean-value inequality spreads the mass of a large
point value over a two-dimensional neighbourhood, not over the line
through the point: from $|\zeta(\sigma_0+it_n)|$ large one can only
conclude that the area integral over a disk is large, hence that
\emph{some} abscissa $\sigma''_n$ (depending on $n$) carries a large
one-dimensional window integral---and there is no way, within this
argument, to pin the moving $\sigma''_n$ to the fixed line
$\sigma_0$. Nor do the classical convexity tools apply: the
logarithmic convexity in $\sigma$ of $L^p$ means of holomorphic
functions on vertical lines is a statement about integrals over the
whole line, not over unit windows, and there is no one-dimensional
sub-mean-value inequality for $|g|^p$, $g$ holomorphic, along a
vertical segment. The band formulation of Theorem~\ref{thm:main} is
not a technical convenience but the exact boundary of what the
subharmonicity method proves.

\subsection{Relation to the Riemann Hypothesis}
The result is a negative theorem about seminorms and is logically
independent of the Riemann Hypothesis (it is unconditional). Its
contextual relevance is this: by a classical theme going back to
Bohr \cite{Bohr22} and made precise by Bagchi \cite{Bagchi}, RH is
equivalent to a strong recurrence (self-approximation) property of
$\zeta$ in compact subsets of the strip, so one might hope to approach
such recurrence statements through an almost-periodicity seminorm
strong enough to control compact windows. Theorem~\ref{thm:main} shows
that the Stepanov scale---the only classical intermediate scale with
that property---is unconditionally unavailable in the band sense; any
window-localised recurrence for $\zeta$ must therefore be obtained by
mechanisms other than membership of $\zeta$ in a classical
almost-periodicity class that sees windows. We make no claim beyond
this.

\section*{Acknowledgements}
The author thanks the broader analytic number theory literature cited
here; all results used are classical or published, and all statements
proved here are elementary consequences of them.

\begin{thebibliography}{9}

\bibitem{Aistleitner}
C.~Aistleitner, \emph{Lower bounds for the maximum of the Riemann zeta
function along vertical lines}, Math. Ann. \textbf{365} (2016),
no.~1--2, 473--496; arXiv:1409.6035.

\bibitem{Bagchi}
B.~Bagchi, \emph{Recurrence in topological dynamics and the Riemann
hypothesis}, Acta Math. Hungar. \textbf{50} (1987), 227--240.

\bibitem{Besicovitch}
A.~S.~Besicovitch, \emph{Almost Periodic Functions}, Cambridge
University Press, Cambridge, 1932.

\bibitem{Bohr22}
H.~Bohr, \emph{\"Uber eine quasi-periodische Eigenschaft Dirichletscher
Reihen mit Anwendung auf die Dirichletschen $L$-Funktionen},
Math. Ann. \textbf{85} (1922), 115--122.

\bibitem{LevitanZhikov}
B.~M.~Levitan and V.~V.~Zhikov, \emph{Almost Periodic Functions and
Differential Equations}, Cambridge University Press, Cambridge, 1982.

\bibitem{Montgomery}
H.~L.~Montgomery, \emph{Extreme values of the Riemann zeta function},
Comment. Math. Helv. \textbf{52} (1977), 511--518.

\bibitem{Titchmarsh}
E.~C.~Titchmarsh, \emph{The Theory of the Riemann Zeta-Function},
2nd ed., revised by D.~R.~Heath-Brown, Oxford University Press,
Oxford, 1986. Chapters VII (mean-value theorems) and VIII
($\Omega$-theorems).

\end{thebibliography}

\end{document}
